\section{The origin section in the line-bundle reduction}

We now specify the local geometric hypotheses under which the
four-torus filling is obtained from \eqref{eq:weierstrass}.
Let $O$ be its origin section, let $P_0$ be a section disjoint
from $O$ near $t=0$, and put $\mathcal M=\mathcal O(P_0-O)$.
Write $\mathcal M^\times$ for the complement of its zero section.
Suppose a specified lifted translation acts on $\mathcal M^\times$
and its quotient over the punctured disk is a family of complex
two-tori. Choose the origin of these tori using a nonzero section
$e$ of $\mathcal M|_O$ over the disk.

After $t=s^3$, let $P'_0,O'$ be the transformed sections on
\eqref{eq:good-equation}, and put
$\mathcal M'=\mathcal O(P'_0-O')$. We assume the following smooth
reduction data: $P'_0$ is disjoint from $O'$ near zero; the specified
lifted translation extends to $(\mathcal M')^\times$ and commutes
with the natural deck action; its quotient is a proper smooth
family $\mathcal F$ of complex two-tori over the $s$ disk. The
period marking is \eqref{eq:period-lattice}, with $\delta$ the
positive circle of the multiplicative fibre and $e_4$ the stated
second vertical period. The positive monodromy is
\eqref{eq:positive-T}.

These are hypotheses on the local smooth reduction and its lifted
translation. The calculation below constructs their marked filling
and attachment maps; it does not establish the existence of that
reduction or of the lifted translation from an arbitrary elliptic
family. In particular it does not replace the needed commutation
or properness by a lattice calculation.

\begin{lemma}\label{lem:conormal-comparison}
There is a nonzero section $e'$ of $\mathcal M'|_{O'}$ such that
the canonical divisor identification over the punctured disk sends
the pulled-back section $e$ to $s^{-2}e'$. The natural deck action
sends $e'(s)$ to $\zeta e'(\zeta s)$.
\end{lemma}
\begin{proof}
Since $P_0$ does not meet $O$ in this neighbourhood,
$\mathcal O(P_0-O)$ is canonically the ideal line bundle of $O$
there. Restricting this ideal to $O$ gives its conormal line:
the ideal modulo its square. Similarly $\mathcal M'|_{O'}$ is
the conormal line of $O'$, because $P'_0$ is disjoint from it.

The local parameters transverse to the origin sections are
$w=-x/y$ and $W=-X/Y$. To check this at infinity, use the
projective chart with $Y\ne0$: the cubic equation solves $Z/Y$
as a holomorphic function of $X/Y$ near the origin, so $-X/Y$
is a local parameter. The Weierstrass substitution gives the
exact equality
\begin{equation}\label{eq:conormal-coordinate}
                         w=s^{-2}W.
\end{equation}
The nonzero section $e$ can therefore be written as $u(t)w$ in
the conormal line, for a holomorphic unit $u$ on a sufficiently
small disk. Define $e'(s)=u(s^3)W$. This is holomorphic and
nonzero, and the canonical divisor identification on $s\ne0$
gives $e=s^{-2}e'$. No choice of a comparison multiplied by an
extra power of $s$ is made here: the parameters are identified
by the actual birational coordinate substitution.

The original pulled-back section is equivariant for the deck
action. If $\rho e'(s)=\lambda(s)e'(\zeta s)$, applying $\rho$
to $e=s^{-2}e'$ consequently yields
$\lambda(s)=(\zeta s)^{-2}/s^{-2}=\zeta$.
The equality extends to zero. Equivalently, the forward action on
the conormal line is the inverse of the normal-coordinate factor
$\zeta^2$ in \eqref{eq:actual-elliptic-deck}.
\end{proof}

Use the image of $e'$ as the origin of the smooth torus family
$\mathcal F$. An automorphism between compact complex tori is
affine in these origins: a lift to their complex vector spaces
has lattice-periodic derivative, whose entries are holomorphic
functions on a compact torus and hence constant. Subtracting its
linear part leaves a constant translation. Its linear part is
determined by its map on periods. Lemma~\ref{lem:inverse-descent}
therefore gives linear part $A=T^{-1}$, while
Lemma~\ref{lem:conormal-comparison} gives translation $\delta/3$.
Indeed multiplication by $e^{2\pi i/3}$ in the multiplicative
fibre is precisely one third of its positive period $\delta$.

Period coordinates make these maps constant over the smooth disk:
their integral linear matrices are locally constant, and the
translation is the specified torsion section $\delta(s)/3$.
Consequently the original deck transformation and its twist by
the unchanged $a_0=e_4/3$ are, respectively,
\begin{equation}\label{eq:actual-affine-actions}
 (z,s)\longmapsto(Az+\delta/3,\zeta s),\qquad
 (z,s)\longmapsto(Az+(\delta+e_4)/3,\zeta s).
\end{equation}
Thus the central affine term does not vanish in the origin $e'$.
The second action is free: in the original $e_4$ coordinate its
first two powers translate by $1/3$ and $2/3$.

Along a positive root $s=e^{2\pi iv/3}$ the original origin $e$
differs from $e'$ by $-2v\delta/3$. The logarithmic twist adds
$ve_4/3$. A small nonunit boundary radius can be absorbed by a
constant rescaling of the chosen frame; the angular formula is
unchanged. At the initial point $s=1$ the two origins coincide.
This determines the marked boundary path, including its original
section, and not only the homology monodromy.

\section{The marked local filling and its integral chains}

Write $\delta=U=u_0+u_1+u_2$ and retain the original coordinates
$(x,t)$ in the basis $(u_0,u_1,u_2,e_4)$. Under the preceding
smooth reduction hypotheses, the twisted local filling is
\begin{equation}\label{eq:marked-native-quotient}
 N_0=(X\times S^1_t\times D)/
 \bigl((x,t,s)\sim(Rx+U/3,t+1/3,\zeta s)\bigr).
\end{equation}
Its boundary is marked by the original section $e$ and the
counterclockwise meridian. The formula just derived is
\begin{equation}\label{eq:marked-native-j}
 j_0([x,t,v])=[x-2vU/3,t+v/3,e^{2\pi iv/3}].
\end{equation}
At $v=1$, the inverse of the generator in
\eqref{eq:marked-native-quotient} sends the endpoint to
$[Px-U,t,1]=[Px,t,1]$. Thus its return is $P$, as required in
the original period marking; $qj_0=e^{2\pi iv}$.

\begin{theorem}\label{thm:marked-native-filling}
The filling \eqref{eq:marked-native-quotient} has the following
explicit orientation-preserving disk trivialization:
\begin{equation}\label{eq:marked-native-Phi}
 \Phi_0[x,t,s]=([x+2tU,3t],e^{-2\pi it}s)
                 \in G_P\times D.
\end{equation}
Here $q=e^{2\pi i\tau}z^3$. The marked fibre and positive
boundary maps followed by retraction are
\begin{equation}\label{eq:marked-native-FK}
 f_0(x,t)=[x+2tU,3t],\qquad
 k_0([x,t,v])=[x+2tU,3t+v].
\end{equation}
The integral fibre basis
\begin{equation}\label{eq:marked-cell-basis}
             u_0,u_1,u_2,v_4,\qquad v_4=e_4-2\delta
\end{equation}
has determinant one relative to $(e_1,e_2,e_3,e_4)$ and is
preserved by $T$, which acts as $P\times I$. In its actual
product cells the filling complex and maps are
\begin{equation}\label{eq:marked-final-chains}
 \begin{aligned}
 Q_{0,n}&=D_n\oplus D_{n-1},&d_0(x,y)&=((P-I)y,0),\\
 F_{0,n}(x,y)&=(x,Sy),&H_{0,n}(x,y)&=(0,x),\\
 K_{0,n}(x,y,z,w)&=(x,Sy+z).&&
 \end{aligned}
\end{equation}
In particular $d_0H_0+H_0d_C=F_0(M-I)$ with the positive
native monodromy and the actual marked boundary section.
\end{theorem}
\begin{proof}
Set $y=x+2tU$. The change $(x,t)\mapsto(y,t)$ is an integral
torus diffeomorphism of determinant one, with inverse
$x=y-2tU$. Under the generator of \eqref{eq:marked-native-quotient},
\[
 y\longmapsto Rx+U/3+2(t+1/3)U=Ry+U=Ry
                         \quad\text{in }X.
\]
It therefore identifies this quotient with
\eqref{eq:positive-quotient} in the $y$ coordinates, without
changing $s$ or the orientation of its disk. It also sends the
marked boundary path to
$[y,t+v/3,e^{2\pi iv/3}]$, since the terms $-2vU/3$ and
$2(v/3)U$ cancel. Applying Theorem~\ref{thm:positive-geometry}
gives \eqref{eq:marked-native-Phi} and
\eqref{eq:marked-native-FK}, with their orientations and inverses.

In the original coordinates, the $t$ loop with $y=0$ has vector
$e_4-2U$. Hence its product-cell basis is exactly
\eqref{eq:marked-cell-basis}. Its matrix of columns is
\begin{equation}\label{eq:marked-basis-matrix}
 V=\begin{pmatrix}
 0&-1&-1&4\\0&0&-1&2\\1&1&1&-6\\0&0&0&1
 \end{pmatrix}.
\end{equation}
It is obtained from the previously verified determinant-one
permutation basis by adding $-2(u_0+u_1+u_2)$ to the last
column, so its determinant is one. Since $T$ fixes both
$e_4$ and $U$, it fixes $v_4$ and permutes the other cells.
These are therefore genuine transported product cells, obtained
by an integral diffeomorphism, in every degree.

In these cells the maps in \eqref{eq:marked-native-FK} are
$[y,3t]$ and $[y,3t+v]$. Their interval subdivisions are exactly
those in Theorem~\ref{thm:positive-chains}, which proves all
formulas in \eqref{eq:marked-final-chains} and all their chain
identities. Proposition~\ref{prop:positive-splitting} gives the
section and kernel of the full boundary map. In particular the
degree-four map and all adjacent prism terms are included.
\end{proof}

The original twist has not been changed: in this cell basis
$e_4/3=(v_4+2U)/3$. Its sum with the original central translation
$U/3$ is $v_4/3+U$, which equals $v_4/3$ as a torus translation.
Both contributions are needed for the integral comparison.

For a homology class written in the original period basis, the
same maps are explicit as well. In degree $n$, change its exterior
coordinates by $\bigwedge^nV^{-1}$, express the $u$-pairs in
$(w_0,w_1,w_2)$, and move $v_4$ to the front in the suspended
terms before applying \eqref{eq:marked-final-chains}. Moving it
past $n-1$ fibre vectors contributes $(-1)^{n-1}$.
Thus this is an integral change of actual geometric
cells, with no rational replacement of the lattice.
For example, write $c=[u_0]=[u_1]=[u_2]$ and let $h$ be the
positive core interval in $H_1(G_P;\mathbb Z)$. Then
\begin{equation}\label{eq:marked-degree-one}
 f_{0*}(e_1)=f_{0*}(e_2)=0,\qquad
 f_{0*}(e_3)=c,\qquad f_{0*}(e_4)=6c+3h,
 \qquad k_{0*}(\gamma_0)=h.
\end{equation}
Indeed $e_1=u_0-u_1$, $e_2=u_1-u_2$, $e_3=u_0$,
and $e_4=v_4+2U$. This gives a deciding distinction from the
zero-origin descent in the unchanged original marking.
At the level of the fundamental group, three traversals of
the marked meridian satisfy
\begin{equation}\label{eq:marked-meridian-relation}
                   \gamma_0^3=f_{0\#}(e_4-2\delta).
\end{equation}
To verify it directly, lift \eqref{eq:marked-native-j} to the
universal vector coordinates. Its endpoint is the affine deck
generator followed by translation by $-\delta$. Cubing adds
$e_4+\delta-3\delta=e_4-2\delta$. Equivalently in $G_P$, the
marked loop is its positive core interval, while the $v_4$ fibre
loop covers three such intervals. The image of $\delta$ is
$3c$ and is not zero in this local filling. Consequently the
local relation cannot be replaced by $\gamma_0^3=f_{0\#}(e_4)$.

For the higher degrees, abbreviate $e_i\wedge e_j$ by $e_{ij}$
and similarly for three or four indices. In the ordered domain
bases $(e_{12},e_{13},e_{14},e_{23},e_{24},e_{34})$ and
$(e_{123},e_{124},e_{134},e_{234})$, the degree-two and
degree-three maps are, respectively,
\begin{equation}\label{eq:marked-higher-matrices}
 f_{0*,2}=\begin{pmatrix}3&1&0&-2&0&0\\0&0&0&0&0&-1\end{pmatrix},
 \qquad
 f_{0*,3}=\begin{pmatrix}1&6&2&-4\\0&3&1&-2\end{pmatrix}.
\end{equation}
The target bases are $([w_0],\sigma U)$ and
$(\omega,\sigma W)$ in the core complex. For example,
$e_{12}=w_0+w_1+w_2$ gives the first coefficient $3$.
The term $e_3\wedge v_4=-v_4\wedge u_0$ gives the last
coefficient $-1$ in degree two. In degree three,
$e_{124}=v_4\wedge e_{12}+2e_{12}\wedge U$
gives $3\sigma W+6\omega$. Expanding the other columns
in the same stated basis gives \eqref{eq:marked-higher-matrices}.
Both maps are integrally surjective: $e_{13},-e_{34}$ map to
the two target generators in degree two, and
$e_{123},e_{134}-2e_{123}$ do so in degree three.
Their kernels have the following full integral bases:
\[
 \begin{array}{c|l}
 2&e_{12}-3e_{13},\quad e_{14},\quad 2e_{13}+e_{23},\quad e_{24},\\
 3&e_{124}-3e_{134},\quad 2e_{134}+e_{234}.
 \end{array}
\]
Indeed the two equations in each displayed matrix solve for the
remaining coordinates with integer coefficients. In degree four,
$e_{1234}=-v_4\wedge\omega$, so its image is $-3\sigma\omega$.
The covering has geometric degree $+3$: the interval-first target
cell $\sigma\omega$ is the negative of the fibre-first core
orientation. These formulas retain both the signs and the complete
integral images in the original periods.

All homology groups, kernel ranks, and cokernels in
Theorem~\ref{thm:homology} remain valid in the stated cell basis.
The invariant degree-one core covector pulls back to
$e_3^*+6e_4^*$, and the core-circle covector pulls back to
$3e_4^*$. Their image is still
$\mathbb Ze_3^*\oplus3\mathbb Ze_4^*$. The higher-degree
specialization indices remain $1,1,3$ in degrees two, three,
and four, by the explicit unimodular change of cells.
Equal indices therefore do not imply equal marked attachment maps.

This local conclusion depends on the smooth line-bundle reduction,
its commuting lifted translation, and its stated period marking.
Within those hypotheses the conormal comparison constructs the
affine deck action and the positive boundary map. Comparisons with
a common chain complex carrying the other local monodromies still
require explicit maps and homotopies. The other finite-order filling,
the cusp attachments, and their global integral gluing have not been
constructed here. No conclusion about the homology of their union
or a complex structure on a sphere follows from this local calculation.
